\documentclass[10pt]{article}
\usepackage[margin=20mm]{geometry}
\usepackage{amsmath,amssymb,mathtools,fontspec,unicode-math,xurl,hyperref}
\IfFontExistsTF{DejaVu Serif}{\setmainfont{DejaVu Serif}}{\setmainfont{Latin Modern Roman}}
\setmathfont{latinmodern-math.otf}
\setlength{\parindent}{0pt}\setlength{\parskip}{6pt}\sloppy
\title{Initial review of
A new proof of Chen’s theorem for Markoff graphs}
\author{}
\date{26 October 2026}
\begin{document}
\maketitle
Abstract: The acquired publisher version of Daniel E. Martin’s four-page paper has a correct involution-summing argument for \(p>3\). Its unrestricted divisibility statement needs the exceptional prime \(3\) removed: for the displayed Markoff equation the graph at \(p=3\) is an eight-vertex cube. This is a local scope correction. The large-prime connectivity consequence survives. The full proof, zero-coordinate cases, involution fixed points and external group-action convention were checked.\par
\section*{Identity and outcome}
M25-0085 is the publisher version of record, Inventiones mathematicae 241 (2025), 623–626, DOI 10.1007/s00222-025-01346-9 [paper]. Printed page numbers are PDF page numbers plus 622. All four original pages were actually viewed, and all 224 extracted text lines were read. The manifest contains the publisher PDF and publisher-page evidence, without author code or TeX; their absence is not counted as a defect.\par
The outcome is Major: the central divisibility claim needs an exceptional-prime restriction. There is a complete counterexample to the all-prime wording of Theorem 1.4. There is no defect established in its intended \(p>3\) argument, and no reason from this finding to reject its all-but-finitely-many-primes consequence.\par
\section*{Definitions and complete verification for \(p>3\)}
The vertex set is \[V_p=\{(x_1,x_2,x_3)\in\mathbb F_p^3\setminus\{0\}:
 x_1^2+x_2^2+x_3^2=3x_1x_2x_3\}.\] Edges are furnished by the three individual Vieta involutions \(R_i\), with \(R_1(x)=(3x_2x_3-x_1,x_2,x_3)\). Coordinate permutations and quotienting by sign changes are not part of this graph definition. Each \(R_i\) preserves the equation, is its own inverse, and cannot send a nonzero vertex to zero.\par
A nonzero vertex cannot have two zero coordinates: the equation would force the remaining square, hence that coordinate, to vanish. Therefore the zero-coordinate branches in (2.1), printed page 624, are disjoint and cover all cases where the product is zero. For nonzero product define \(y_i=x_i/(3x_jx_k)\). If \(x_i=0\), put \(y_i=0\) and the other two values equal to \(1/2\). All denominators used are invertible when \(p>3\). Dividing the Markoff equation by \(3x_1x_2x_3\), or using the zero-coordinate definition, gives \[\sum_{i=1}^3 y_i(x)=1.\]\par
For each \(i\), the identity \[y_i(x)+y_i(R_i x)=1\] holds in all cases. If the other two coordinates are nonzero, it is the root-sum identity divided by \(3x_jx_k\), including the case \(x_i=0\). If one of the other coordinates is zero, both values are \(1/2\). In particular, a fixed point of \(R_i\) has \(y_i=1/2\). Such points must not be counted as two distinct vertices.\par
For a connected component \(C\), \(R_i\) is a bijection of \(C\). Consequently the argument can be written without any ambiguous edge multiplicity: \[2\sum_{x\in C}y_i(x)
 =\sum_{x\in C}\bigl(y_i(x)+y_i(R_i x)\bigr)=|C|
 \quad\hbox{in }\mathbb F_p.\] Summing in \(i\) yields \(2|C|=3|C|\), hence \(|C|=0\) in \(\mathbb F_p\), equivalently \(p\mid |C|\). This proves the main argument for every prime \(p>3\), including components containing zeros or involution fixed points.\par
\section*{Problem: Theorem 1.4 needs an exception at \(p=3\)}
\subsection*{Exact location and counterexample}
Theorem 1.4, printed page 624 (PDF page 2), states divisibility without excluding a prime. The abstract has the same unrestricted wording. The proof on that page explicitly begins with \(p>3\).\par
Set \(p=3\). The right-hand side of the defining equation is zero. A nonzero square in \(\mathbb F_3\) equals one. Thus a solution has either zero or three nonzero coordinates; excluding the zero triple leaves exactly \[V_3=\{1,-1\}^3.\] Each \(R_i\) replaces \(x_i\) by \(-x_i\) and leaves the other coordinates fixed. These independent sign flips connect every one of the eight vertices to every other. Hence the entire graph is one component, isomorphic to a three-dimensional cube: \[|C|=8,\qquad 3\nmid8.\] This is an exact finite classification under the paper’s own equation and graph definition, not an inference from a numerical pattern.\par
For comparison, \(p=2\) gives the four vertices \((1,1,1),(0,1,1),(1,0,1),(1,1,0)\). The central vertex \((1,1,1)\) is connected to each of the other three, so the graph is connected and has size four, divisible by two. Its loops do not change the component. Thus the complete corrected statement for this equation is divisibility for every prime \(p\ne3\).\par
\subsection*{Counterevidence, impact and repair limits}
The explicit \(p>3\) opening of the proof is strong evidence that the author intended the restriction. It does not make the unrestricted theorem wording true. Add \(p>3\) to Theorem 1.4 and the corresponding abstract statement; alternatively state \(p\ne3\) and append the elementary \(p=2\) check.\par
The equation with coefficient \(1\), obtained by rescaling coordinates by \(3\), is equivalent to this one only when \(p\ne3\). It cannot be used to dismiss this counterexample. Likewise, connectivity itself is not contradicted: \(G_3\) is connected. The failure is its asserted component divisibility.\par
This correction does not alter the proof for \(p>3\), the large-prime result, or any effective bound claimed only for sufficiently large primes. The error is Major for the stated central divisibility guarantee: Theorem 1.4 and the abstract assert it without a prime restriction, and the eight-vertex component at \(p=3\) satisfies the stated graph definition while violating that guarantee. The failure is confined to this domain boundary, and the explicit restriction above repairs it. Repairability does not remove the counterexample to the published central statement. No author adoption of the repair is claimed.\par
\section*{Dependency and group-convention check}
The all-but-finitely-many-primes inference uses the Bourgain–Gamburd–Sarnak large-orbit theorem. Its primary text [bgs], pages 1–2, defines \(\Gamma\) using both the Vieta involutions and coordinate permutations, and gives a large \(\Gamma\)-orbit with complement at most \(p^\varepsilon\). The proof of that external theorem was not re-proved here.\par
This convention difference can be bridged after the divisibility argument, as follows. Let \(H=\langle R_1,R_2,R_3\rangle\). Permutations normalize \(H\) and fix \(v=(1,1,1)\), so \(\Gamma v=Hv\). Choose, for example, \(\varepsilon=1/2\). The \(H\)-orbit of \(v\) has size at least \(p\); it cannot lie outside the large \(\Gamma\)-orbit, whose total complement is less than \(p\). Therefore that large orbit equals \(Hv\). Any other \(H\)-orbit would lie in the complement and have size both at least \(p\) and less than \(p\), a contradiction. This supplies the convention bridge without assuming arbitrary \(\Gamma\)-orbits are single \(H\)-orbits, and without using the connectivity conclusion to prove divisibility.\par
The separately cited numerical threshold \(3.489\cdot10^{392}\) was not independently reconstructed from the external effective estimates. It is not needed for the elementary divisibility proof. No assertion is made here that all exceptional small primes have been verified beyond the finite range below.\par
\section*{Supporting checks}
An independently written Python program enumerated all triples and all Vieta edges for the fifteen primes \[2,3,5,7,11,13,17,19,23,29,31,37,41,43,47.\] It checked equation preservation, involutivity, component sizes, the \(y\)-sum identity, the edge identity, and the three component coordinate sums when \(p>3\). No permutation edges or sign quotients were introduced. Selected outputs are: \[\begin{array}{c|rrrrrr}
p&2&3&5&7&11&13\\ \hline
|V_p|&4&8&40&28&88&208\\
\text{component sizes}&4&8&40&28&88&208\\
\#\{x:R_1x=x\}&2&0&0&4&8&8
\end{array}\] All tested primes had one component. The \(p=3\) divisibility test failed exactly as proved above; all tested \(p>3\) identities passed. These computations corroborate the complete symbolic argument; finite testing is not used to prove the all-prime result.\par
The official arXiv record inspected on 26 October 2026 lists only v1 dated 21 February 2025. That record is contextual; no author version is substituted for the acquired publisher PDF. A focused correction search was nonexhaustive and did not establish an adopted correction.\par
The evidence directory contains the independently written enumeration source and exact JSON output. No author code was executed. The final report is really compiled from its exact current TeX with authorized Tectonic 0.17.0. Pre/post source, compiler and full task-local dependency hashes, actual logs, PDF hash, and actual inspection of every final PDF page are recorded and linked through provenance and delivery state. `human\_reviews=[]`; `scientific\_decision=pending`. Delivery is not scientific acceptance.\par
Focused independent acceptance is identified by token:\par
The current executor used a bounded read-only native helper for the existing problem, with actual roles and evidence recorded in provenance. The retained initial review is distinguished from this focused check; the scientific decision remains pending.\par
9 D. E. Martin, A new proof of Chen’s theorem for Markoff graphs, Invent. Math. 241 (2025), 623–626. https://doi.org/10.1007/s00222-025-01346-9. Assigned PDF SHA-256: `3e1f4b7f99f18b71159463aaf2d9e649e11317272782f9105299a9d72369097c`. J. Bourgain, A. Gamburd and P. Sarnak, Markoff surfaces and strong approximation: 1, arXiv:1607.01530v1, 6 July 2016, pp. 1–2. https://arxiv.org/pdf/1607.01530. D. E. Martin, official arXiv record, https://arxiv.org/abs/2502.15960, accessed 26 October 2026.\par
\end{document}
